\section{Finite heads, boundary residuals and compatible approximation}
\label{sec:finite}

An infinite hierarchy is useful for approximation only when the error
introduced at its boundary can be transported to the observables of
interest. We now make that boundary explicit. This also separates two
notions of finite approximation: truncating the list of moments, and
evolving a finite divisor whose moments remain mutually compatible.
The distinction was motivated by the earlier hierarchy discussion
\cite{BakerCM15}; the calculations below are consequences of
Theorem~\ref{thm:hierarchy}.

\subsection{The four boundary rows}

Freeze the scale $\ell$ from \eqref{eq:domain} and put
\begin{equation}
 \sigma=(t-t_0)/\ell^2,\qquad \delta=D/\ell^2,\qquad
 s_n=\ell^{2n}S_n,\qquad v_n=\ell^{2n-1}T_n,\qquad m=s_1.
 \label{fin:normalization}
\end{equation}
The equations \eqref{eq:S-hierarchy}--\eqref{eq:T-hierarchy} have the
same coefficients after replacing $(D,S,T,t)$ by $(\delta,s,v,\sigma)$.
Fix $P\ge3$ and retain
$x_P=(\delta,s_1,\ldots,s_P,v_2,\ldots,v_P)$.
Define the polynomial vector field $f_P$ by setting all indices above
$P$ to zero in those equations. This defines an approximation, not an
identity for the actual higher moments.

\begin{proposition}[Boundary residual]\label{fin:residual}
The actual canonical head satisfies $x_P'=f_P(x_P)+r_P$, where primes
in this section denote $\sigma$ derivatives. Only four residual rows
can be nonzero:
\begin{align}
 r_{s,P-1}&=-4(P-1)(P+2)\delta s_{P+1},\notag\\
 r_{s,P}&=-4P(P+2)s_{P+1}-4P(P+3)\delta s_{P+2},\notag\\
 r_{v,P-1}&=-4(P-1)(P+2)\delta v_{P+1},\notag\\
 r_{v,P}&=-4(P^2+P-1)v_{P+1}-4P(P+3)\delta v_{P+2}.
 \label{fin:four-rows}
\end{align}
With $\gamma=q_0/a^2<1$ and $B_1=14\Cu\gamma$,
\begin{equation}
 \|r_P\|_\infty\le
 4P(P+3)\sqrt{q_0}B_1L\gamma^{P-1/2}
             \max\{d,1+d\gamma\}.
 \label{fin:residual-bound}
\end{equation}
In particular the residual tends to zero geometrically at fixed height
and collar parameters.
\end{proposition}

\begin{proof}
Every nonlinear factor in row $n$ has index at most $n$. The only
indices larger than $P$ therefore occur in the two linear higher-moment
terms of rows $P-1$ and $P$, giving \eqref{fin:four-rows}.
The shell bound used below in \eqref{eq:integer-bound} gives
$\mathcal A_1=\sum_j|\mu_j|\le B_1L$; it follows directly from
\eqref{eq:count-safe} and \eqref{eq:collar-kernel}. Thus
\[
 |s_n|\le B_1L\gamma^{n-1},\qquad
 |v_n|\le\sqrt{q_0}B_1L\gamma^{n-3/2}\quad(n\ge2).
\]
For the odd estimate use
$|w_j|/\sqrt{|Q_j|}\le\sqrt{q_0}$ and sum
$|\mu_j|^{n-1/2}$. Apply these inequalities to the four rows,
using $|\delta|\le d$, $P^2+P-1\le P(P+3)$ and $\gamma<1$.
The $P-1$ rows contribute the separate factor $d$ in
\eqref{fin:residual-bound}; they have no extra factor $\gamma$.
\end{proof}

\subsection{What an independent absolute error allowance loses}

Suppose $y'=f_P(y)$ and a convex box contains both $x_P$ and $y$
throughout $[0,\Sigma]$. If $K_P$ bounds the induced max-norm Jacobian
of $f_P$ on that box and $\epsilon_P$ is the right side of
\eqref{fin:residual-bound}, the usual differential inequality gives
\begin{equation}
 \|x_P(\tau)-y(\tau)\|_\infty
 \le \e^{K_P\tau}\|x_P(0)-y(0)\|_\infty
       +\epsilon_P\frac{\e^{K_P\tau}-1}{K_P}.
 \label{fin:gronwall}
\end{equation}
At $K_P=0$ the quotient means $\tau$. The containing box is a
premise of this bound. In the present unweighted coordinates,
\[
 \left|\frac{\partial (f_P)_{s,P-1}}{\partial s_P}\right|
       =4(P^2-1),
\]
so a scalar Jacobian majorant grows at least quadratically with $P$.
Combining that particular allowance with a geometrically small boundary
residual does not prove convergence on a fixed positive duration.
A componentwise version makes the source of this loss more precise.

\begin{proposition}[A scoped positive-comparison limitation]
\label{fin:positive-box}
Let $A_P$ be a constant entrywise nonnegative absolute-Jacobian majorant
on the containing box. Give its residual majorant an independent constant
component
\[
 \bar r_{s,P}=4P(P+2)B_1L\gamma^P
             +4P(P+3)dB_1L\gamma^{P+1},
\]
with all other components nonnegative. The comparison solution with
zero initial error,
$z(\tau)=\int_0^\tau\exp((\tau-u)A_P)\bar r\dd u$, satisfies
\begin{equation}
 z_m(\tau)\ge\frac{B_1L}{2}(P+2)(P+1)!(4\gamma\tau)^P.
 \label{fin:path-lower}
\end{equation}
For fixed $\tau>0$, $L$ and $\gamma>0$, this allowance diverges as
$P\to\infty$. The conclusion concerns this independent constant-tail,
absolute-Jacobian comparison, not the signed truncation error.
\end{proposition}

\begin{proof}
Row $s_n$ has derivative $-4n(n+2)$ with respect to $s_{n+1}$;
its nonlinear terms do not contain that component. Hence the
successive entries of $A_P$ along the path
$s_P\to s_{P-1}\to\cdots\to s_1$ are at least $4n(n+2)$.
Their product is
\[
 \prod_{n=1}^{P-1}4n(n+2)
       =4^{P-1}(P-1)!(P+1)!/2.
\]
Every matrix power and every component of $\bar r$ is nonnegative.
Keep this one path in the $(P-1)$st term of the exponential series
and the first term of $\bar r_{s,P}$. Integrating
$(\tau-u)^{P-1}/(P-1)!$ gives $\tau^P/P!$, proving
\eqref{fin:path-lower}. The ratio of successive factorial-exponential
terms tends to infinity, proving the stated divergence.
\end{proof}

The discarded signs and compatibility relations could matter to the
actual solution. Proposition~\ref{fin:positive-box} neither proves
instability of a moment truncation nor rules out a compatible closure.
It explains why a small local residual is not by itself a useful
fixed-duration error theorem. Shorter steps also need a controlled
restart and accumulation argument.

\subsection{Preserving the divisor and the intended error consumer}

A different finite object is obtained by truncating the complete
opposite-paired product at an initial time $t_0$, retaining conjugates,
ties and multiplicities, and evolving the resulting polynomial $P_N$
by the heat equation. Its higher moments are then determined by one
actual divisor; they are not set independently to zero. At the purely
polynomial level the evolution has the exact expression
\begin{equation}
 P_N(t_0+\tau,z)
 =\exp(-\tau\partial_z^2)P_N(t_0,z)
 =\mathbb E\,P_N(t_0,z+\ii\sqrt{2\tau}\,X),\qquad \tau\ge0,
 \label{fin:polynomial-heat}
\end{equation}
where $X$ is standard normal. To verify the second identity, expand the
polynomial in its finite Taylor series. Odd Gaussian moments vanish,
and $\mathbb EX^{2k}=(2k)!/(2^kk!)$ converts the even terms exactly
to $(-\tau)^k\partial_z^{2k}P_N/k!$.

The source in \eqref{eq:source} gives the same Gaussian formula for $H$:
expectation multiplies $\cos(zu)$ by $\exp(\tau u^2)$. Its
super-exponential decay justifies Fubini for bounded real $\tau\ge0$
and bounded complex $z$. What remains unproved here is a common-time
error estimate for $P_N-H$ and the transfer to the selected pair and
moments. The slice-jet passage used in Theorem~\ref{thm:hierarchy}
reinitializes at each reference time and does not supply that estimate.
A possible route is a common product growth envelope, the complete
omitted inverse-square tail $O(\log N/N)$, Gaussian domination, and
Cauchy and Rouch\'e estimates on the independently cleared tube.
This is a specified approximation problem, not a theorem of the present
paper. The cutoff $N$ is a divisor cutoff, distinct from the moment
order $P$; initial canonical divisor data and its tail must be supplied.

The error needed by transport is simple to state. Let $m_N$ be real,
as for a conjugation-compatible approximation. If it satisfies $|m-m_N|\le\eta$ on a scaled interval of length $\Sigma$, then
for every subinterval $[x,y]$,
\begin{equation}
 \left|\log\frac{W(x,y)}{W_N(x,y)}\right|
 \le4(y-x)\eta\le4\Sigma\eta,
 \qquad W(x,y)=\exp\left(-4\int_x^y m\dd\sigma\right).
 \label{fin:transport-error}
\end{equation}
This follows by subtracting the two integrals. The same uniform
subinterval control bounds the inhomogeneous cost multiplicatively;
the physical cost retains its factor $\ell^2$. Forcing error must also
retain the product in $12s_2+16\delta s_3$. None of these consequences
supplies the independent tube or gap anchors.

The analogy with BBGKY truncation is therefore organizational: it asks
how unresolved information enters a retained hierarchy. The older
Jacobi-chain memory variable was observation height, whereas the present
variable is Newman time; the omitted moments here generally have nonzero
initial data. Inter-level compatibility does not by itself imply stability
or physical realizability, and an exact projection-memory formalism needs
a specified state space and flow \cite{AkbariEtAl2012Hierarchy,ChorinHaldKupferman2000}.
No density-matrix positivity, independence law or arbitrary-sequence
well-posedness is assumed for these divisor moments.
